% ============================================================
\section{变量与运算符}
\subsection*{Variables and Operators}

本节介绍 PXL 中的变量与运算符概念。

\subsection{变量}
\subsubsection*{Variables}

变量用于在程序中临时存储信息。
声明变量时，变量名后须跟冒号（\cmd{:}）与变量类型。声明时也可给出初值。见下一小节示例。

\begin{noteBox}
有关命名限制，见「标识符」小节。
\end{noteBox}

\subsubsection{初始化变量}
\paragraph*{Initializing Variables}

可用两种方式将变量初始化为任意受支持的数据类型：
\begin{itemize}
  \item 使用字面量值。
\begin{lstlisting}[style=pxl]
y : boolean = true;                   // initialize by literal value
\end{lstlisting}
  \item 使用表达式。
\begin{lstlisting}[style=pxl]
z : boolean = !y;                     // initialized by expression
\end{lstlisting}
\end{itemize}

下列示例展示不同的初始化方法：
\begin{lstlisting}[style=pxl]
x : integer;                              // no initial value
pi : double = 3.141592;                   // initialized by value
pi2 : double = pi * 2;                    // initialized by expression
s : string = "abc" "def";                 // concatenated
s1 : string = "abc" + pi;                 // converted to string,then concatenated
c : constraint of double = (10.0,200];    // initialized by value
\end{lstlisting}

\subsubsection{变量作用域}
\paragraph*{Scope of Variables}

PXL 中的作用域（scope）指变量与表达式在程序若干上下文中的含义。下文通过示例说明与变量作用域相关的规则与条件。

局部作用域由一对花括号（\cmd{\{\}}）定义。也就是说，用花括号定义的复合语句会引入新的变量作用域。花括号的用法见后文「复合语句」（Compound Statement）一节。

流程控制语句会创建新作用域。
\begin{lstlisting}[style=pxl]
if ( var == 1) {                // Start of new scope
   my_var: integer = 1;
   ...
}                             // End of scope.
\end{lstlisting}

\begin{itemize}
  \item 在局部作用域中声明的变量在该作用域外不存在。也就是说，若变量在某一作用域内声明，则不能在该作用域外使用。例如：
\begin{lstlisting}[style=pxl]
{
   x = y + b;
}
z = x;      // x is not defined; it only exists in the local scope
\end{lstlisting}
  \item 变量从最外层作用域到最内层作用域全局可见。例如：
\begin{lstlisting}[style=pxl]
{
    x : integer = 1;
       {
       a = x + b; // x exists and can be seen within this scope
                   // (global from outer scope)
       x = y + b; // global value of x is overwritten
       }
    z = x;         // x is defined here and contains the value of y + b
}
z = x;                // x is not defined here - parser error
\end{lstlisting}
  \item 在外层作用域声明、又在内层作用域重新声明的变量会被遮蔽（shadowed），以避免值冲突。
  \begin{itemize}
    \item 若发生变量遮蔽，且被遮蔽变量的作用域并非由函数声明定义，则程序生成警告。
    \item 在函数调用内，遮蔽是预期行为，不会警告。例如：
\begin{lstlisting}[style=pxl]
y: integer = 1;
x: integer = 1;
{
   y : integer = 3;      // The original y value is shadowed.
                         // This is a new y.
    x = x + 1;           // x is equal to 2 here
    y = y + 1;           // y is equal to 4 here
}
                         // y is equal to 1 here - the shadowed value
                         // returns when exiting the inner scope
                         // x is equal to 2 here
\end{lstlisting}
  \end{itemize}
\end{itemize}

下面是变量遮蔽的另一示例。函数声明：
\begin{lstlisting}[style=pxl]
fxy : function (
   x : integer;
   y : integer;
) returning z : integer {
      x = x + 1;
      y = y + 1;
      z = x + y;
};
\end{lstlisting}

函数调用：
\begin{lstlisting}[style=pxl]
x1 : integer = 1;
y : integer = 1;

z1 = fxy(x1, y);                              // within the function y is shadowed
                                              // y in the local scope is "2"
x2 = x1;                                      // x1 is "1"
y1 = y;                                       // y is still "1"
z2 = z1;                                      // z1 is "4"
\end{lstlisting}

\subsection{运算符}
\subsubsection*{Operators}

运算符是表达式中用于指定在求值时应执行某些操作的符号。

PXL 支持下列运算符：
\begin{itemize}
  \item 赋值运算符（Assignment Operators）
  \item 关系运算符（Relational Operators）
  \item 逻辑运算符（Logical Operators）
  \item 条件运算符（Conditional Operator）
  \item 数值与按位运算符（Numeric and Bitwise Operators）
  \item 仅输出 double 的运算符（Operators With Double-Only Outputs）
\end{itemize}

\subsubsection{赋值运算符}
\paragraph*{Assignment Operators}

赋值运算符对表达式求值并将其保存到目标左值（lvalue）。
左值是可容纳值、并可出现在赋值语句左侧的表达式。

涉及在约束值中使用 \cmd{>} 或 \cmd{>=} 的赋值，赋值运算符前后必须有空白。也可将约束值放在赋值 \cmd{=} 之后的括号内。

例如：
\begin{lstlisting}[style=pxl]
abc = >=3;
xyz = (>5);
\end{lstlisting}

\begin{noteBox}
赋值运算符不可串联。下例不合法：\\
\cmd{a = b = c}
\end{noteBox}

正确赋值示例如下：
\begin{lstlisting}[style=pxl]
x:      integer;      // declare a variable named x

point: newtype struct of { x: integer; y: integer; }; // declare a struct
                                                      // of integers
p: point;
x   = 42;      // assign x a value of 42
p.x = 10;      // assign p.x a value of 10
p.y = 27;      // assign p.y a value of 27
\end{lstlisting}

\subsubsection{关系运算符}
\paragraph*{Relational Operators}

关系表达式将表达式求值为布尔值 \cmd{true} 或 \cmd{false}。表~\ref{tab:pxl-rel-ops} 定义关系运算符。

\begin{longtable}{@{}>{\raggedright\arraybackslash\ttfamily}p{0.16\textwidth} p{0.42\textwidth} p{0.34\textwidth}@{}}
\caption{关系运算符（Relational Operators）}\label{tab:pxl-rel-ops}\\
\toprule
\textbf{\textrm{语法}} & \textbf{\textrm{说明}} & \textbf{\textrm{适用于}} \\
\midrule
\endfirsthead
\toprule
\textbf{\textrm{语法}} & \textbf{\textrm{说明}} & \textbf{\textrm{适用于}} \\
\midrule
\endhead
\bottomrule
\endfoot
e1 == e2 & 若 e1 等于 e2 则为 true，否则为 false & 所有完全定义的数据类型 \\
e1 != e2 & 若 e1 不等于 e2 则为 true，否则为 false & 所有完全定义的数据类型 \\
e1 < e2 & 若 e1 小于 e2 则为 true，否则为 false & 数值 \\
e1 <= e2 & 若 e1 不大于 e2 则为 true，否则为 false & 数值 \\
e1 >= e2 & 若 e1 不小于 e2 则为 true，否则为 false & 数值 \\
e1 > e2 & 若 e1 大于 e2 则为 true，否则为 false & 数值 \\
\end{longtable}

\subsubsection{逻辑运算符}
\paragraph*{Logical Operators}

逻辑表达式求值为布尔值。逻辑运算符要求全部参数均为 \cmd{boolean} 类型。表达式从左到右求值，并在结果已确定时停止。表~\ref{tab:pxl-logic-ops} 定义逻辑运算符。

\begin{longtable}{@{}>{\raggedright\arraybackslash\ttfamily}p{0.22\textwidth} p{0.70\textwidth}@{}}
\caption{逻辑运算符（Logical Operators）}\label{tab:pxl-logic-ops}\\
\toprule
\textbf{\textrm{语法}} & \textbf{\textrm{说明}} \\
\midrule
\endfirsthead
\toprule
\textbf{\textrm{语法}} & \textbf{\textrm{说明}} \\
\midrule
\endhead
\bottomrule
\endfoot
e1 \&\& e2 & 若 e1 为 false，则结果为 false，否则为 e2 \\
e1 || e2 & 若 e1 为 true，则结果为 true，否则为 e2 \\
! e1 & 若 e1 为 true，则结果为 false，否则为 true \\
\end{longtable}

\subsubsection{条件运算符}
\paragraph*{Conditional Operator}

条件表达式接受三个参数；第一个参数必须为 \cmd{boolean} 类型。表~\ref{tab:pxl-cond-op} 定义条件运算符。

\begin{longtable}{@{}>{\raggedright\arraybackslash\ttfamily}p{0.22\textwidth} p{0.70\textwidth}@{}}
\caption{条件运算符（Conditional Operator）}\label{tab:pxl-cond-op}\\
\toprule
\textbf{\textrm{语法}} & \textbf{\textrm{说明}} \\
\midrule
\endfirsthead
\toprule
\textbf{\textrm{语法}} & \textbf{\textrm{说明}} \\
\midrule
\endhead
\bottomrule
\endfoot
e1 ? e2 : e3 & 若 e1 为 true，则为 e2，否则为 e3 \\
\end{longtable}

考虑带有三个参数 \cmd{e1}、\cmd{e2}、\cmd{e3} 的条件运算符。求值顺序如下：
\begin{enumerate}
  \item 求值第一个参数（\cmd{e1}）。
  \item 若第一个参数（\cmd{e1}）为 \cmd{true}，则求值第二个参数。
  \item 若第一个参数（\cmd{e1}）为 \cmd{false}，则求值第三个参数（\cmd{e3}）。
\end{enumerate}

\subsubsection{数值与按位运算符}
\paragraph*{Numeric and Bitwise Operators}

PXL 支持全部标准数学运算符。按位运算符仅支持整数。表~\ref{tab:pxl-num-ops} 定义数值与按位运算符。

\begin{longtable}{@{}>{\raggedright\arraybackslash\ttfamily}p{0.14\textwidth} p{0.22\textwidth} p{0.24\textwidth} p{0.32\textwidth}@{}}
\caption{数值与按位运算符（Numeric and Bitwise Operators）}\label{tab:pxl-num-ops}\\
\toprule
\textbf{\textrm{语法}} & \textbf{\textrm{说明}} & \textbf{\textrm{允许类型}} & \textbf{\textrm{注释}} \\
\midrule
\endfirsthead
\toprule
\textbf{\textrm{语法}} & \textbf{\textrm{说明}} & \textbf{\textrm{允许类型}} & \textbf{\textrm{注释}} \\
\midrule
\endhead
\bottomrule
\endfoot
e1 + e2 & 加法 & Integer 与 double & 任一参数为 double 时返回 double \\
e1 - e2 & 减法 & Integer 与 double & 任一参数为 double 时返回 double \\
e1 * e2 & 乘法 & Integer 与 double & 任一参数为 double 时返回 double \\
e1 \% e2 & 整数余数 & Integer 与 double & 任一参数为 double 时返回 double \\
e1 << e2 & 逻辑左移（移入零） & Integer &  \\
e1 >> e2 & 逻辑右移（移入零） & Integer &  \\
e1 \& e2 & 按位 AND & Integer &  \\
e1 | e2 & 按位 OR & Integer &  \\
e1 \^{} e2 & 按位 XOR & Integer &  \\
! e1 & 按位 NOT & Integer &  \\
\end{longtable}

\subsubsection{仅输出 double 的运算符}
\paragraph*{Operators With Double-Only Outputs}

除法（\cmd{/}）与幂运算（\cmd{**}）将参数视为 double，并且始终只返回 double。表~\ref{tab:pxl-double-ops} 定义这些运算符。

\begin{longtable}{@{}>{\raggedright\arraybackslash\ttfamily}p{0.14\textwidth} p{0.20\textwidth} p{0.24\textwidth} p{0.34\textwidth}@{}}
\caption{仅输出 double 的运算符（Operators With Double-Only Outputs）}\label{tab:pxl-double-ops}\\
\toprule
\textbf{\textrm{语法}} & \textbf{\textrm{说明}} & \textbf{\textrm{接受}} & \textbf{\textrm{注释}} \\
\midrule
\endfirsthead
\toprule
\textbf{\textrm{语法}} & \textbf{\textrm{说明}} & \textbf{\textrm{接受}} & \textbf{\textrm{注释}} \\
\midrule
\endhead
\bottomrule
\endfoot
e1 / e2 & 除法 & Integer 与 double & e1 与 e2 被强制为 double；运算始终返回 double \\
e1 ** e2 & 幂运算 & Double & e1 与 e2 被强制为 double；运算始终返回 double \\
\end{longtable}

\begin{noteBox}
任何含除法运算符（\cmd{/}）或幂运算符（\cmd{**}）的表达式，均不可用作列表或哈希表的索引。
\end{noteBox}

若要执行整数运算，必须将该运算包在 \cmd{dtoi()} 函数调用中。该函数将 double 转换为 integer。例如：
\begin{lstlisting}[style=pxl]
var : integer = dtoi(10/3);
\end{lstlisting}

\subsection{成员引用}
\subsubsection*{Member Reference}

成员引用表达式访问复合对象的一个成员。成员引用取该单一成员的类型。该成员是左值，并可作为赋值目标。表~\ref{tab:pxl-member-ref} 定义成员引用表达式。
有关左值的更多信息，见「赋值运算符」。

\begin{longtable}{@{}>{\raggedright\arraybackslash\ttfamily}p{0.22\textwidth} p{0.70\textwidth}@{}}
\caption{成员引用（Member Reference）}\label{tab:pxl-member-ref}\\
\toprule
\textbf{\textrm{语法}} & \textbf{\textrm{说明}} \\
\midrule
\endfirsthead
\toprule
\textbf{\textrm{语法}} & \textbf{\textrm{说明}} \\
\midrule
\endhead
\bottomrule
\endfoot
e1 . e2 & 按名称访问结构体成员 \\
e1 [ e2 ] & 按索引访问列表或哈希成员 \\
\end{longtable}

\subsection{优先级与结合性}
\subsubsection*{Precedence and Associativity}

表~\ref{tab:pxl-precedence} 列出 PXL 中全部运算符的优先级与结合性。表顶优先级最高，表底最低。

\begin{longtable}{@{}>{\raggedright\arraybackslash\ttfamily}p{0.48\textwidth} p{0.44\textwidth}@{}}
\caption{运算符的优先级与结合性（最高优先级在前）}\label{tab:pxl-precedence}\\
\toprule
\textbf{\textrm{运算符}} & \textbf{\textrm{结合性}} \\
\midrule
\endfirsthead
\toprule
\textbf{\textrm{运算符}} & \textbf{\textrm{结合性}} \\
\midrule
\endhead
\bottomrule
\endfoot
() \quad [] \quad . & 从左到右 \\
** & 从右到左（幂运算） \\
! \quad + \quad - & 从右到左（一元运算符） \\
== \quad != \quad < \quad <= \quad > \quad >= & 非法（一元运算符） \\
* \quad / \quad \% & 从左到右 \\
+ \quad - & 从左到右 \\
<< \quad >> & 从左到右 \\
< \quad <= \quad > \quad >= & 非结合 \\
== \quad != & 非结合 \\
\& & 从左到右 \\
\^{} & 从左到右 \\
| & 从左到右 \\
\&\& & 从左到右 \\
|| & 从左到右 \\
?: & 从右到左 \\
= \quad @= \quad ?= & 非法 \\
\textrm{用户定义运算符} & 非结合 \\
\end{longtable}

表~\ref{tab:pxl-associativity} 定义各类结合性。

\begin{longtable}{@{}p{0.28\textwidth} p{0.64\textwidth}@{}}
\caption{结合性（Associativity）}\label{tab:pxl-associativity}\\
\toprule
\textbf{结合性} & \textbf{定义} \\
\midrule
\endfirsthead
\toprule
\textbf{结合性} & \textbf{定义} \\
\midrule
\endhead
\bottomrule
\endfoot
从左到右（Left to right） & 表达式从左到右求值。 \\
从右到左（Right to left） & 表达式从右到左求值。 \\
非结合（Nonassociative） & 需要括号。 \\
非法（Illegal） & 该运算符不可串联。运算符产生的类型与其消费的类型不同。 \\
\end{longtable}

下列示例说明结合性如何工作：
\begin{lstlisting}[style=pxl]
/*
 * pairs of expressions show evaluation order
 */

A + B + C
( A + B ) + C

A ** B ** C
A ** ( B ** C )

A * B + C >> 2
( ( A * B ) + C ) >> 2

A >= B == C
( A >= B ) == C

A < B == C > D
( A < B ) == ( C > D )
\end{lstlisting}
